Das Verfassen schriftlicher Ausarbeitungen ist faktisch in allen Studienfächern
unumgänglich. Neben einem fachlich fundierten Inhalt trägt auch eine
übersichtliche Gestaltung des Textbildes wesentlich zur Verständlichkeit einer
Arbeit bei. Diese Vorlage entstand, um die Hürden für die Benutzung des
Textsatzsystems \LaTeX\ zur Erstellung studentischer Arbeiten möglichst
niedrig zu legen. Aufgrund einer Vielzahl von unterschiedlichen Vorgaben, ist
diese Vorlage vielmehr als Baukasten und Sammlung von Beispielen zu verstehen
und erhebt nicht den Anspruch einer Layout-Vorgabe. Eigene Kommandos, werden nur
an den Stellen definiert, wo es unumgänglich ist (z.\,B. bei der Gestaltung der
Deckblätter etc.). Hierdurch ist es auch dem fortgeschrittenen
\LaTeX-Anwender möglich, eigene Anpassungen in diese Vorlage einzupflegen.

Kapitel~\ref{sec:vorlage} befasst sich mit dem Aufbau und der Funktionsweise
dieser Vorlage. Außerdem wird hier erläutert wie sich die Vorlage für den
jeweiligen Einsatzzweck einrichten lässt. Diese Kapitel sollte daher
keinesfalls übersprungen werden, da es dabei hilft, sich schnell im Grundgerüst
dieser Vorlage zurecht zu finden.

In Kapitel~\ref{sec:latex} wird eine Sammlung von \LaTeX-Beispielen für
verschiedene Anwendungsfälle bereitgehalten. Die entsprechenden Unterabschnitte
sind weitestgehend unabhängig voneinander und sollten nur bei Bedarf gelesen
werden.

Wenn die Fertigstellung des Textkörpers vorangeschritten ist, kann bereits die
formale Überarbeitung der Arbeit beginnen. Kapitel~\ref{sec:final} fasst
hierfür wichtige Hinweise zu diesem Thema zusammen. Bei der Erstellung eines
Deckblatts und der verschiedenen Verzeichnisse (Inhalts-, Abbildungs-,
Literaturverzeichnis, etc.) nimmt \LaTeX\ Ihnen viel Arbeit ab. Der Aufwand
für die Kontrolle auf korrekte Rechtschreibung und inhaltliche Stimmigkeit darf
aber nicht unterschätzt werden. Planen Sie mindestens zwei Wochen zeitlichen
Puffer ein. Gerade bei umfangereicheren Arbeiten (Bachelorarbeit bzw.
Diplomarbeit und aufwärts) erfolgt der Druck in einer Druckerei. Für einen
reibungslosen Ablauf sind einige Fallstricke zu beachten. Auch hier können Sie
mit einen kleinen zeitlichen Zusatzpuffer von einer halben bis einer Woche für
die Endkontrolle der Druckvorlage viel Stress vermeiden. Überspringen
Sie dieses Kapitel also \textbf{nicht}, wenn es um's Ganze geht.


% https://www.tu-chemnitz.de/mathematik/part_dgl/teaching/material/Hinweise_fuer_Abschlussarbeiten.pdf
% http://osg.informatik.tu-chemnitz.de/lehre/richtlinien/ausarbeitungen.pdf

% Zielgruppe
% - Studenten
%   -> Projekt, Seminar-, Abschlussarbeiten
% - Promovenden, Habilitanten
%   -> Dissertationen, Habil.
%   -> Buchveröffentlichungen
